\begin{figure}[htbp]
			\centering
			\resizebox{\textwidth}{!}{%
				\begin{tikzpicture}[
					scale=1.05,
					asta/.style={line width=1.5pt, line cap=round, line join=round, black},
					puntone/.style={line width=3.5pt, line cap=round, line join=round, black},
					tirante/.style={line width=1.0pt, line cap=round, line join=round, black},
					nodo/.style={circle, fill=white, draw=black,
						inner sep=0pt, minimum size=6pt, line width=1.2pt},
					forza/.style={->, >=stealth, line width=1.5pt, black},
					sez/.style={black, thick, dashed},
					]
					\def\SEP{8.5}
					
					%%──────── PANNELLO (a): stato di sollecitazione ────────
					\begin{scope}
						\coordinate (A) at (0,0); \coordinate (B) at (3,0); \coordinate (C) at (6,0);
						\coordinate (D) at (2,1); \coordinate (E) at (4,1);
						\draw[puntone](A)--(D);
						\draw[puntone](D)--(E);
						\draw[puntone](B)--(E);
						\draw[puntone](E)--(C);
						\draw[tirante](A)--(B);
						\draw[tirante](B)--(C);
						\draw[tirante](D)--(B);
						\draw[line width=1.2pt](A)--++(-0.4,-0.65)--++(0.8,0)--cycle;
						\draw[line width=1.2pt](-0.55,-0.65)--(0.55,-0.65);
						\foreach \x in {-0.45,-0.25,-0.05,0.15,0.35}\draw[thin](\x,-0.65)--++(-.18,-.22);
						\draw[line width=1.2pt](C)--++(-0.4,-0.50)--++(0.8,0)--cycle;
						\draw[line width=1.2pt]($(C)+(-0.50,-0.50)$)--($(C)+(0.50,-0.50)$);
						\draw[line width=1.2pt]($(C)+(-0.25,-0.65)$)circle(3.5pt);
						\draw[line width=1.2pt]($(C)+( 0.25,-0.65)$)circle(3.5pt);
						\draw[line width=1.2pt]($(C)+(-0.55,-0.82)$)--($(C)+(0.55,-0.82)$);
						\foreach \x in {-0.45,-0.25,-0.05,0.15,0.35}\draw[thin]($(C)+(\x,-0.82)$)--++(-.18,-.22);
						\node[nodo]at(A){};\node[nodo]at(B){};\node[nodo]at(C){};\node[nodo]at(D){};\node[nodo]at(E){};
						\draw[forza]($(D)+(-1.25,0)$)--($(D)+(-0.12,0)$);
						\node[left,font=\small]at($(D)+(-1.25,0)$){$2F$};
						\draw[forza]($(E)+(0,1.20)$)--($(E)+(0,0.12)$);
						\node[above,font=\small]at($(E)+(0,1.20)$){$3F$};
						\node[above left=-1pt,font=\small]at(A){$A$};
						\node[below=3pt,font=\small]at(B){$B$};
						\node[above right=-1pt,font=\small]at(C){$C$};
						\node[above=2pt,font=\small]at(D){$D$};
						\node[above right=0pt,font=\small]at(E){$E$};
						\node[font=\normalsize]at(3,-1.6){$(a)$};
					\end{scope}
					
					%%──────── PANNELLO (b): sezione di Ritter ────────
					\begin{scope}[xshift=\SEP cm]
						\coordinate (A) at (0,0); \coordinate (B) at (3,0); \coordinate (C) at (6,0);
						\coordinate (D) at (2,1); \coordinate (E) at (4,1);
						\draw[asta](A)--(B);\draw[asta](B)--(C);\draw[asta](A)--(D);
						\draw[asta](D)--(B);\draw[asta](D)--(E);\draw[asta](B)--(E);\draw[asta](E)--(C);
						\draw[sez](2.4,-0.95)--(2.4,1.75) node[above,font=\small]{sezione};
						\draw[line width=1.2pt](A)--++(-0.4,-0.65)--++(0.8,0)--cycle;
						\draw[line width=1.2pt](-0.55,-0.65)--(0.55,-0.65);
						\foreach \x in {-0.45,-0.25,-0.05,0.15,0.35}\draw[thin](\x,-0.65)--++(-.18,-.22);
						\draw[line width=1.2pt](C)--++(-0.4,-0.50)--++(0.8,0)--cycle;
						\draw[line width=1.2pt]($(C)+(-0.50,-0.50)$)--($(C)+(0.50,-0.50)$);
						\draw[line width=1.2pt]($(C)+(-0.25,-0.65)$)circle(3.5pt);
						\draw[line width=1.2pt]($(C)+( 0.25,-0.65)$)circle(3.5pt);
						\draw[line width=1.2pt]($(C)+(-0.55,-0.82)$)--($(C)+(0.55,-0.82)$);
						\foreach \x in {-0.45,-0.25,-0.05,0.15,0.35}\draw[thin]($(C)+(\x,-0.82)$)--++(-.18,-.22);
						\node[nodo]at(A){};\node[nodo]at(B){};\node[nodo]at(C){};\node[nodo]at(D){};\node[nodo]at(E){};
						\draw[forza]($(D)+(-1.25,0)$)--($(D)+(-0.12,0)$);
						\node[left,font=\small]at($(D)+(-1.25,0)$){$2F$};
						\draw[forza]($(E)+(0,1.20)$)--($(E)+(0,0.12)$);
						\node[above,font=\small]at($(E)+(0,1.20)$){$3F$};
						\node[above=2pt,font=\small,gray!70]at(D){$P_{AB}$};
						\node[below=3pt,font=\small,gray!70]at(B){$P_{DE}$};
						\node[above left=-1pt,font=\small]at(A){$A$};
						\node[above right=-1pt,font=\small]at(C){$C$};
						\node[below=2pt,font=\small]at(D){$D$};
						\node[above right=0pt,font=\small]at(E){$E$};
						\node[font=\normalsize]at(3,-1.6){$(b)$};
					\end{scope}
					
				\end{tikzpicture}
			}
			\caption{$(a)$~Stato di sollecitazione: tratto
				spesso le aste compresse (puntoni $AD$, $DE$,
				$BE$, $EC$), tratto sottile le aste tese
				(tiranti $AB$, $BC$, $DB$).
				$(b)$~Sezione di Ritter sul pannello centrale:
				il taglio attraversa $AB$, $DB$ ed $DE$. I poli
				dei due correnti coincidono con nodi della
				travatura, $P_{DE}\equiv B$ e $P_{AB}\equiv D$;
				poiché i correnti $AB$ e $DE$ sono paralleli,
				il polo della diagonale $DB$ è all'infinito e
				$N_{DB}$ si ricava dall'equilibrio alla
				traslazione verticale.}
			\label{fig:tre_maglie_soluzione}
		\end{figure}
